\chapter{Pre- and post-processing}
\label{chap:python}

\framepart{Outline}

\cindex{post-processing, in Python}\sindex{"@excel1}%
\cindex{Python}\cindex{matplotlib}\cindex{pandas}\cindex{pyvista}%
A finite element study is rarely one computation. It is a parameter varied
over a range, a mesh refined until the answer stops moving, a model compared
with an experiment -- and then all of it done again when something changes.
Doing that by editing a \texttt{.dgibi} file by hand and copying numbers into
a spreadsheet works for the first three runs and becomes unmanageable at the
twentieth.

This chapter is about the two ends of that loop: generating the input and
analysing the output, from outside \texttt{Cast3M}. It gives you, in order:

\begin{itemize}
\item a modern replacement for \op{@excel1} on the \texttt{Cast3M} side, which
  writes any number of curves to one CSV file with a header
  (section~\ref{sec:py:export});
\item how to read that file and plot it (section~\ref{sec:py:curves});
\item how to read a whole field -- mesh, displacement, stress -- and
  post-process it (section~\ref{sec:py:fields});
\item how to drive a parametric study from Python
  (section~\ref{sec:py:param});
\item for those who would rather stay in the shell, how to do the same
  with a shell script, \texttt{sed} and \texttt{awk}, in \texttt{bash} or in
  \texttt{csh} (section~\ref{sec:py:shell});
\item how to wrap \texttt{Cast3M} as a function of its parameters and hand
  it to OpenTURNS or JAX (section~\ref{sec:py:drivers});
\item and where to get these tools on Linux, macOS and Windows
  (section~\ref{sec:py:install}).
\end{itemize}

\noindent The last of these is the \emph{pre}-processing half, and it is
worth reading even if you intend to use Python: the shell approach is the one
that survives on a cluster, in a batch queue, and on any machine where you
cannot install a Python package.

\noindent You need \texttt{numpy}, \texttt{pandas}, \texttt{matplotlib} for
the curves and \texttt{pyvista} (or \texttt{meshio}) for the fields:

\begin{codebox}
\begin{verbatim}
$ pip install numpy pandas matplotlib pyvista meshio h5py
\end{verbatim}
\end{codebox}

\noindent \texttt{h5py} is needed only to read MED files, which are HDF5
underneath.

\section{Getting the data out: a better \texttt{@excel1}}
\label{sec:py:export}

\op{@excel1} takes one \op{evolution} and writes two columns. In practice you
almost always have several curves sharing the same abscissa -- a time, a
loading parameter -- and you want them in one file, with a header line naming
them, so that the Python side needs no manual editing.

\sindex{EXPORTCSV (procedure)}\cindex{CSV export}%
\dindex{index}{list the indices of a table}%
\dindex{dime[nsion]}{size of a table or a list}%
The following procedure does that. It takes a \op{table} of \op{evolution}
objects and writes a semicolon-separated file whose first line holds the
column names.

\begin{castemcom}
The structure is the one of \op{@excel1}: save the current output unit,
redirect to the file, print, restore.

The differences are that the loop over the table indices builds the header
line first, that all the ordinates are collected into a table of
\op{listreel} before printing, and that the abscissa is taken from the first
evolution.

\op{dime} gives the number of entries of a table or the length of a list, and
\op{index} returns the list of indices of a table, which is what makes the
procedure work whatever you have put in it. The header and each data line are
assembled as strings with \op{mot} and \op{chai} and glued together with
\op{et}.

\medskip
Save it as \texttt{procedur/exportcsv.procedur} next to your
\texttt{.dgibi} file and it will be found automatically -- see the end of
chapter~\ref{chap:importexport}.
\end{castemcom}
\begin{castemlines}
\begin{verbatim}
DEBPROC EXPORTCSV TEVOL*TABLE
                  FICOUT*MOT;

  LIDX = INDEX TEVOL;
  NC   = DIME LIDX;

  EV1  = TEVOL . (EXTR LIDX 1);
  ABS1 = EXTR EV1 ABSC;
  NP   = DIME ABS1;

  II = VALE IMPR;
  OPTI IMPR 10 IMPR FICOUT;

*  -- header line --
  LIGN = MOT 'ABSC';
  REPETER BC NC;
    IC   = EXTR LIDX &BC;
    LIGN = LIGN ET (MOT ';')
                ET (MOT IC);
  FIN BC;
  MESSAGE LIGN;

*  -- one line per point --
  REPETER BP NP;
    XV   = EXTR ABS1 &BP;
    LIGN = CHAI XV;
    REPETER BC NC;
      EVC  = TEVOL . (EXTR LIDX &BC);
      ORDC = EXTR EVC ORDO;
      YV   = EXTR ORDC &BP;
      LIGN = LIGN ET (MOT ';')
                  ET (CHAI YV);
    FIN BC;
    MESSAGE LIGN;
  FIN BP;

  OPTI IMPR II;
FINPROC;
\end{verbatim}
\end{castemlines}

\begin{castemcom}
Using it is a one-liner. Collect your curves in a table, index them by the
name you want to see as a column header, and call the procedure.

This is the natural end of the time loop of chapter~\ref{chap:heat} or of the
\op{pasapas} post-processing of chapter~\ref{chap:plasticity}.
\end{castemcom}
\begin{castemlines}
\begin{verbatim}
TEV = TABLE;
TEV . 'UY_A3' = evol manu
   'TEMPS' thetime 'UY' uuy;
TEV . 'SMYY'  = evol manu
   'TEMPS' thetime 'SMYY' sigyy;
TEV . 'SMXY'  = evol manu
   'TEMPS' thetime 'SMXY' sigxy;

EXPORTCSV TEV 'history.csv';
\end{verbatim}
\end{castemlines}

\noindent The result is an ordinary CSV file:

\begin{codebox}
\begin{verbatim}
ABSC;UY_A3;SMYY;SMXY
0;0;0;0
0.5;7.8e-05;7845.91;2353.77
1;0.000156;15643.4;4693.03
...
\end{verbatim}
\end{codebox}

\noindent For fields rather than curves, do not use CSV at all: write a VTK
file with \op{sort 'VTK'}, as in section~\ref{sec:io:meshexchange}. That is
what section~\ref{sec:py:fields} reads.

\section{Curves: pandas and matplotlib}
\label{sec:py:curves}

Reading the file and drawing a publication-quality figure takes a dozen
lines. The header line means the columns are addressed by name, so the script
does not have to change when you add a curve.

\begin{codebox}
\begin{verbatim}
import pandas as pd
import matplotlib.pyplot as plt

df = pd.read_csv("history.csv", sep=";")
print(df.head())          # always look at what you actually read

fig, ax = plt.subplots(figsize=(6, 4))
ax.plot(df["ABSC"], df["SMYY"] / 1e6, "-o", ms=3, label=r"$\sigma_{yy}$")
ax.plot(df["ABSC"], df["SMXY"] / 1e6, "-s", ms=3, label=r"$\sigma_{xy}$")
ax.set_xlabel("time")
ax.set_ylabel("stress [MPa]")
ax.legend()
ax.grid(alpha=0.3)
fig.tight_layout()
fig.savefig("history.pdf")      # vector output, drops straight into LaTeX
\end{verbatim}
\end{codebox}

\noindent Three things this buys you over the spreadsheet route:

\begin{itemize}
\item \texttt{savefig("history.pdf")} produces a vector figure you include
  with \verb|\includegraphics| -- no screenshots, no resolution loss;
\item the script is reproducible: rerun the computation, rerun the script,
  the figure updates. Nobody has to remember which columns were selected in
  which spreadsheet;
\item everything pandas can do is now available -- resampling, filtering,
  comparing several runs:

\begin{codebox}
\begin{verbatim}
runs = {n: pd.read_csv(f"history_{n}.csv", sep=";") for n in (5, 10, 20, 40)}
for n, d in runs.items():
    ax.plot(d["ABSC"], d["SMYY"], label=f"nel = {n}")
\end{verbatim}
\end{codebox}
\end{itemize}

\section{Fields: pyvista}
\label{sec:py:fields}

\cindex{VTK format}\cindex{Paraview}\cindex{field, plotting outside Cast3M}%
For anything defined over the mesh, export a VTK file from \texttt{Cast3M}
and read it with \texttt{pyvista}. This is where Python goes well beyond what
the built-in viewer can do.

\begin{codebox}
\begin{verbatim}
import numpy as np
import pyvista as pv

mesh = pv.read("result.vtu")
print(mesh)                      # cells, points, bounds
print(mesh.array_names)          # ['U', 'SMISES', ...]

vm = mesh["SMISES"]
print(f"max von Mises = {vm.max()/1e6:.1f} MPa at node {vm.argmax()}")
\end{verbatim}
\end{codebox}

\paragraph{Deformed shape with a field on it.} The classic picture, in four
lines. \texttt{factor} is the displacement magnification:

\begin{codebox}
\begin{verbatim}
warped = mesh.warp_by_vector("U", factor=200.0)

p = pv.Plotter(window_size=(900, 700))
p.add_mesh(warped, scalars="SMISES", cmap="viridis",
           scalar_bar_args={"title": "von Mises [Pa]"})
p.add_mesh(mesh, style="wireframe", color="grey", opacity=0.3)
p.view_xy()
p.show()                         # or p.screenshot("vmises.png")
\end{verbatim}
\end{codebox}

\noindent On a machine with no display -- a cluster, a CI job -- add
\texttt{pv.start\_xvfb()} at the top and use
\texttt{pv.Plotter(off\_screen=True)}.

\paragraph{A field along a path.} This is the Python counterpart of
\op{evol 'CHPO'}, and it is strictly more flexible: the sampling points are
yours to choose and need not be mesh nodes, so two meshes can be compared on
the same abscissa.

\begin{codebox}
\begin{verbatim}
line = mesh.sample_over_line((0.2, 0.0, 0.0), (1.0, 0.0, 0.0), resolution=200)

import matplotlib.pyplot as plt
plt.plot(line["Distance"], line["SMISES"] / 1e6)
plt.xlabel("distance from the hole [m]")
plt.ylabel(r"$\sigma_{vm}$ [MPa]")
\end{verbatim}
\end{codebox}

\paragraph{Integrating over the domain.} Useful to check a resultant, a mean
value or the dissipated energy:

\begin{codebox}
\begin{verbatim}
sized = mesh.compute_cell_sizes(area=True)
area  = sized.cell_data["Area"].sum()
print(f"domain area = {area:.4f}")

# area-weighted mean of a cell field
mean = (sized.cell_data["Area"] * sized.cell_data["SOMEFIELD"]).sum() / area
\end{verbatim}
\end{codebox}

\paragraph{Dropping into pandas.} When you want the raw numbers rather than a
picture:

\begin{codebox}
\begin{verbatim}
import pandas as pd

df = pd.DataFrame(mesh.points[:, :2], columns=["x", "y"])
df["smises"]      = mesh["SMISES"]
df[["ux", "uy"]]  = mesh["U"][:, :2]

print(df.sort_values("smises", ascending=False).head())   # hot spots
df.to_csv("nodal_results.csv", index=False)
\end{verbatim}
\end{codebox}

\noindent One caveat on the names: \op{sort 'VTK' mod u sig;} exports the
stress \emph{tensor}, whose components are \texttt{SMXX}, \texttt{SMYY},
\ldots, not a von Mises scalar. If you want \texttt{SMISES} in the file,
compute it in \texttt{Cast3M} first and export it too:

\begin{codebox}
\begin{verbatim}
sigvm = vmis mod sig;
opti sort 'result.vtu';
sort 'VTK' mod u sig sigvm;
\end{verbatim}
\end{codebox}

\noindent or compute it in Python from the tensor components, which is just
as easy and keeps the \texttt{Cast3M} file shorter.

\paragraph{Reading MED directly.} If you prefer to skip the VTK step,
\texttt{meshio} reads the MED files \texttt{Cast3M} writes, and keeps the
group names:

\begin{codebox}
\begin{verbatim}
import meshio

m = meshio.read("plate.med")
print([(c.type, len(c.data)) for c in m.cells])   # [('line',50),('triangle',390)]
print(m.cell_data["cell_tags"])                  # group numbers
print(m.cell_tags)                               # {-1:['BAS'], -5:['TROU'], ...}

meshio.write("plate.vtu", m)                        # convert to anything else
\end{verbatim}
\end{codebox}

\section{Driving a parametric study from Python}
\label{sec:py:param}

\cindex{parametric study}\cindex{reproducible workflow}%
Once reading the results is scripted, it is a short step to generating the
input as well. The pattern is: write the \texttt{.dgibi} file from a
template, run \texttt{castem26} on it, read the result, repeat.

\begin{codebox}
\begin{verbatim}
import subprocess, pathlib
import pandas as pd

TEMPLATE = """
opti dime 2 elem qua4;
ldom = 1.;  r = {radius};
nel  = {nel};
* ... the rest of the computation ...
EXPORTCSV TEV 'out_{tag}.csv';
fin;
"""

results = {}
for r in (0.05, 0.10, 0.20, 0.30):
    tag  = f"r{int(r*100):03d}"
    path = pathlib.Path(f"run_{tag}.dgibi")
    path.write_text(TEMPLATE.format(radius=r, nel=20, tag=tag))

    subprocess.run(["castem26", str(path)], check=True,
                   stdout=subprocess.DEVNULL)

    results[r] = pd.read_csv(f"out_{tag}.csv", sep=";")

# stress concentration factor against hole radius
kt = {r: d["SMYY"].max() / 1e6 for r, d in results.items()}
print(kt)
\end{verbatim}
\end{codebox}

\noindent Two practical notes. \texttt{Cast3M} writes \texttt{fort.3} and
\texttt{fort.98} in the working directory, so \textbf{run each case in its own
directory} (pass \texttt{cwd=} to \texttt{subprocess.run}) if you want to run
several in parallel. And \texttt{castem26} returns a zero exit status even
when the computation fails, so check the output for the string
\texttt{ERREUR} rather than relying on \texttt{check=True} alone.

\section{Running a series of computations from a shell script}
\label{sec:py:shell}

\cindex{shell script}\cindex{batch computation}\cindex{parameter sweep}%
Everything in section~\ref{sec:py:param} can be done without Python, with a
shell script, \texttt{sed} and \texttt{awk}. Three reasons to know how:

\begin{itemize}
\item it needs nothing installed. On a cluster, in a batch queue, on a
  colleague's machine, the shell is already there and your Python environment
  is not;
\item it composes with the job scheduler. A \texttt{for} loop becomes an
  array job with almost no change;
\item for simply substituting a few numbers into a file and pulling a few
  numbers back out, it is shorter than the Python equivalent.
\end{itemize}

\noindent Use Python instead as soon as the \emph{analysis} becomes
structured -- interpolation, fitting, anything with more than one table, or
any plot you intend to publish.

\subsection{\texttt{sed}: putting parameters into a \texttt{.dgibi} file}

\cindex{sed}%
The idea is to write the \texttt{.dgibi} file once, with placeholders where
the parameters go, and let \texttt{sed} replace them. Mark the placeholders
distinctively -- double underscores are a good choice, because they cannot
occur in \textsc{gibiane} by accident:

\begin{codebox}
\begin{verbatim}
* plate.template -- the parametrised computation
opti dime 2 elem qua4;
ldom = 1.;
r    = __RADIUS__;
nel  = __NEL__;
* ... mesh, model, solve, EXPORTCSV ...
fin;
\end{verbatim}
\end{codebox}

\noindent \texttt{sed} then substitutes and the loop runs the cases:

\begin{codebox}
\begin{verbatim}
#!/bin/bash
set -eu                      # stop on the first error

for r in 0.05 0.10 0.20; do
  for nel in 10 20 40; do

    tag="r${r}_n${nel}"
    mkdir -p "run_$tag"

    sed -e "s/__RADIUS__/$r/" \
        -e "s/__NEL__/$nel/"  \
        plate.template > "run_$tag/calcul.dgibi"

    ( cd "run_$tag" && castem26 calcul.dgibi > sortie.log 2>&1 )

    if grep -q 'ERREUR' "run_$tag/sortie.log"; then
      echo "FAILED: $tag"
    fi

  done
done
\end{verbatim}
\end{codebox}

\noindent Four points in that script are worth more than the rest of it:

\begin{itemize}
\item \textbf{each case runs in its own directory.} \texttt{Cast3M} writes
  \texttt{fort.3} and \texttt{fort.98} in the working directory, so two
  simultaneous runs in one directory corrupt each other. The subshell
  \texttt{( cd \ldots \&\& \ldots )} changes directory for the run only and
  leaves you where you started;
\item \textbf{the double quotes around \texttt{"\$tag"}} matter the first time
  a path contains a space;
\item \textbf{\texttt{castem26} returns success even when the computation
  fails}, so the exit status tells you nothing. Grep the output for
  \texttt{ERREUR}, as above;
\item \textbf{keep the template and the generated files separate.} The
  generated \texttt{calcul.dgibi} is disposable; the template is the thing you
  edit and keep under version control.
\end{itemize}

\noindent To run the cases in parallel, append an \verb|&| to the subshell so
it goes into the background,

\begin{codebox}
\begin{verbatim}
    ( cd "run_$tag" && castem26 calcul.dgibi > sortie.log 2>&1 ) &
\end{verbatim}
\end{codebox}

\noindent and add a \texttt{wait} after the loops. This is safe only because
each case has its own directory, which is the whole reason for the
\texttt{mkdir}.

\subsection{\texttt{awk}: processing files with columns}

\cindex{awk}\cindex{gawk}%
\texttt{awk} reads a file line by line, splits each line into fields, and
runs a small program on them. For column data -- which is exactly what
\op{EXPORTCSV} and the \op{@excel1} procedure write -- that is usually all
you need. \texttt{gawk} is the GNU implementation; the examples below are
plain \texttt{awk} and work with either.

The structure of an \texttt{awk} program is \texttt{pattern \{ action \}},
with \texttt{\$1}, \texttt{\$2}, \ldots\ the fields of the current line,
\texttt{NR} the line number and \texttt{NF} the number of fields. The
separator is set with \texttt{-F}.

\begin{codebox}
\begin{verbatim}
# the largest value of column 3, and when it occurred
$ awk -F';' 'NR>1 && $3>m {m=$3; t=$1}
             END {printf "max SMYY = %.4g MPa at t = %s\n", m/1e6, t}' history.csv
max SMYY = 30.99 MPa at t = 2

# convert units and reformat, keeping a header
$ awk -F';' 'NR==1 {print "time  SMYY[MPa]"; next}
             {printf "%-5s %8.3f\n", $1, $3/1e6}' history.csv

# the last line of every run, gathered into one table
$ for d in run_*; do
>   printf "%s " "$d"
>   tail -1 "$d/history.csv"
> done > sweep.dat
\end{verbatim}
\end{codebox}

\noindent \texttt{NR>1} skips the header line; \texttt{next} stops processing
the current line and reads the following one. That last loop is the usual way
a parameter sweep is reduced to a single file ready for plotting.

\subsection{\texttt{csh}: the same loop in the C shell}
\label{sec:py:csh}

\cindex{csh}\cindex{tcsh}\cindex{shell script!C shell}%
Many computing centres, and many older \texttt{Cast3M} installations, still
give their users the C shell (\texttt{csh}, today always \texttt{tcsh}) as
login shell, and batch scripts inherited from colleagues are often written in
it. The loop of the previous pages reads as follows in \texttt{csh}:

\begin{codebox}
\begin{verbatim}
#!/bin/csh -f
# sweep.csh -- the parameter sweep, in C shell
set castem = castem26

foreach r ( 0.05 0.10 0.20 )
  foreach nel ( 10 20 40 )

    set tag = "r${r}_n${nel}"
    mkdir -p "run_$tag"

    sed -e "s/__RADIUS__/$r/" \
        -e "s/__NEL__/$nel/"  \
        plate.template > "run_$tag/calcul.dgibi"

    ( cd "run_$tag" && $castem calcul.dgibi >& sortie.log )

    grep -q 'ERREUR' "run_$tag/sortie.log"
    if ( $status == 0 ) then
      echo "FAILED: $tag"
    endif

  end
end
\end{verbatim}
\end{codebox}

\noindent The logic is identical; the differences are of syntax only:

\begin{itemize}
\item \textbf{loops and tests}: \texttt{foreach \ldots\ end} and
  \texttt{if ( \ldots\ ) then \ldots\ endif} replace \texttt{for \ldots\ done}
  and \texttt{if \ldots\ fi};
\item \textbf{variables} are set with \texttt{set name = value} (spaces
  allowed), and the exit status of the last command is \texttt{\$status}, not
  \texttt{\$?};
\item \textbf{redirection}: \texttt{>\&} sends both the output and the errors
  to the file, where \texttt{bash} writes \texttt{> file 2>\&1};
\item the option \texttt{-f} on the first line stops \texttt{csh} from reading
  your \texttt{.cshrc}, so that the script behaves the same for every user;
\item there is no \texttt{set -e}: test the result of each important command
  yourself, as with \texttt{grep} above.
\end{itemize}

\noindent \texttt{csh} has no functions and an awkward quoting of special
characters, which is why new scripts are better written for \texttt{bash}.
Keep \texttt{csh} for reading and adapting what you inherit.

\subsection*{Which tool for which job}

\begin{gbox}{Choosing between them}
\begin{tabular}{@{}>{\ttfamily}l@{\hspace{5mm}}p{124mm}@{}}
sed    & substitute text in a file. Use it to inject parameters into a
         template, and for nothing else -- it has no arithmetic \\[1.mm]
awk    & column arithmetic and reformatting: extract, convert units, find a
         maximum, reduce many files to one table \\[1.mm]
grep   & find out whether a run failed, or pull one labelled line out of a log \\[1.mm]
Python & everything structured: interpolation, curve fitting, several tables
         at once, fields rather than columns, and any figure for a report \\
\end{tabular}
\end{gbox}

\noindent The honest division of labour is that \texttt{sed} and the shell
generate the runs, \texttt{awk} reduces the output to a small table, and
Python takes over from there. Many working setups use all three, and there
is nothing wrong with that.

\section{\texttt{Cast3M} as a function: OpenTURNS and JAX}
\label{sec:py:drivers}

\cindex{driver, of a computation}\cindex{black box, Cast3M as}%
The parametric study of section~\ref{sec:py:param} runs a list of cases
chosen in advance. An optimizer, a calibration or an uncertainty study
chooses the next case itself, from the results of the previous ones. All it
needs is \texttt{Cast3M} wrapped as a \emph{function} of its parameters. Four
rules make one: a template; one directory per run; a check of the output for
\texttt{ERREUR}, since the exit status says nothing; and a cache, because an
optimizer often asks twice for the same point.

\begin{codebox}
\begin{verbatim}
import pathlib, subprocess, functools
import numpy as np

TEMPLATE = pathlib.Path("plate.template").read_text()

@functools.lru_cache(maxsize=None)          # the same r is computed once
def smax(r):
    """Largest SMYY (MPa) for a hole of radius r: one Cast3M run."""
    wd = pathlib.Path("runs") / f"r{r:.6f}"
    wd.mkdir(parents=True, exist_ok=True)
    (wd / "calcul.dgibi").write_text(
        TEMPLATE.replace("__RADIUS__", f"{r:.6f}").replace("__NEL__", "20"))
    with open(wd / "sortie.log", "w") as log:
        subprocess.run(["castem26", "calcul.dgibi"], cwd=wd,
                       stdout=log, stderr=subprocess.STDOUT)
    if "ERREUR" in (wd / "sortie.log").read_text():
        raise RuntimeError(f"Cast3M failed in {wd}")
    data = np.loadtxt(wd / "history.csv", delimiter=";", skiprows=1)
    return data[:, 1].max() / 1e6
\end{verbatim}
\end{codebox}

\noindent Save it as \texttt{drivers.py}: the two examples below import it.
The template is the one of section~\ref{sec:py:shell}, ending with
\texttt{EXPORTCSV TEV 'history.csv';}.

\subsection{OpenTURNS: propagating an uncertainty}
\label{sec:py:openturns}

\cindex{OpenTURNS}\cindex{uncertainty propagation}%
OpenTURNS, the uncertainty library developed by EDF, Airbus, Phimeca and
IMACS, sees a model as a function from an input vector to an output vector.
\texttt{PythonFunction} turns \texttt{smax} into such a model; a probability
law on the radius then gives a random output, sampled by as many
\texttt{Cast3M} runs as there are points:

\begin{codebox}
\begin{verbatim}
import openturns as ot
from drivers import smax

model = ot.PythonFunction(1, 1, lambda x: [smax(round(x[0], 6))])
R = ot.Normal(0.10, 0.005)                  # radius: mean 0.10, 5% scatter
Y = ot.CompositeRandomVector(model, ot.RandomVector(R))
ot.RandomGenerator.SetSeed(0)
sample = Y.getSample(200)                   # 200 Cast3M runs
print("mean", sample.computeMean()[0],
      "std", sample.computeStandardDeviation()[0],
      "99% quantile", sample.computeQuantile(0.99)[0])
\end{verbatim}
\end{codebox}

\noindent The same model object serves the rest of the library: Sobol'
indices when there are several uncertain parameters, a metamodel (polynomial
chaos, kriging) fitted on a few dozen runs and sampled instead of
\texttt{Cast3M}, or the calibration of material parameters against
measurements (\texttt{NonLinearLeastSquaresCalibration}). For the last one, a
complete example with \texttt{Cast3M} as the model is the identification of
an elastoplastic truss in the lecture notes \emph{Nonlinear Problems in
Mechanics} (Chapter~7 and its \texttt{Cast3M} companion). OpenTURNS also has
its own tools to fill templates and read results,
\texttt{openturns.coupling\_tools}; the few lines of \texttt{smax} do the
same, visibly.

\subsection{JAX: a \texttt{Cast3M} run inside a differentiated program}
\label{sec:py:jax}

\cindex{JAX}\cindex{automatic differentiation}%
JAX differentiates Python programs written with its \texttt{numpy}
(\texttt{jax.numpy}). It cannot look inside \texttt{Cast3M}, but it lets you
\emph{declare} the derivative of a black box: \texttt{jax.pure\_callback}
calls the external code, \texttt{jax.custom\_vjp} says how to multiply a
vector by its derivative. Everything written around the black box -- a cost,
a change of variables, a penalty -- is then differentiated by JAX. Here the
derivative is a central difference of two runs; a derivative computed by
\texttt{Cast3M} itself would replace the function \texttt{\_df} and nothing
else.

\begin{codebox}
\begin{verbatim}
import numpy as np, jax, jax.numpy as jnp
from scipy.optimize import minimize
from drivers import smax
jax.config.update("jax_enable_x64", True)
f64 = jax.ShapeDtypeStruct((), jnp.float64)

def _f(r):                                  # the black box, on numbers
    return np.float64(smax(round(float(r), 6)))

def _df(r, h=1e-3):                         # its derivative: two runs
    r = float(r)
    return np.float64((smax(round(r + h, 6)) - smax(round(r - h, 6))) / (2*h))

@jax.custom_vjp
def s(r):
    return jax.pure_callback(_f, f64, r)

s.defvjp(lambda r: (s(r), r),                               # forward
         lambda r, v: (v * jax.pure_callback(_df, f64, r),)) # backward

target = 2.4                                # wanted max SMYY (MPa)
cost = lambda r: 0.5 * (s(r) - target) ** 2 # any JAX code around s
grad = jax.grad(cost)                       # uses the declared derivative

res = minimize(lambda x: float(cost(x[0])), [0.10],
               jac=lambda x: [float(grad(x[0]))], method="BFGS")
print("radius", res.x[0], "runs", smax.cache_info().currsize)
\end{verbatim}
\end{codebox}

\noindent Check such a gradient before trusting it: compare
$J(r+h)-J(r)$ with $h\,J'(r)$ for decreasing $h$ -- the difference must
decrease like $h^2$ (the Taylor test).

\section{Installing the tools}
\label{sec:py:install}

\cindex{installation, of the tools}%
\texttt{Cast3M} itself is installed as in section~\ref{sec:intro:install}.
The tools of this chapter come from elsewhere:

\begin{gbox}{Where each tool comes from}
\small
\begin{tabular}{@{}>{\raggedright\arraybackslash}p{22mm}@{\hspace{3mm}}>{\raggedright\arraybackslash}p{38mm}@{\hspace{3mm}}>{\raggedright\arraybackslash}p{36mm}@{\hspace{3mm}}>{\raggedright\arraybackslash}p{38mm}@{}}
 & \textbf{GNU/Linux} & \textbf{macOS} & \textbf{Windows} \\[1.mm]
\texttt{bash}, \texttt{sed}, \texttt{awk}
  & installed & installed (\texttt{bash} 3.2; \texttt{zsh} is the login shell)
  & WSL2, or Git Bash \\[1.mm]
\texttt{csh} (\texttt{tcsh})
  & \texttt{apt install tcsh}, \texttt{dnf install tcsh}
  & installed (\texttt{/bin/csh})
  & inside WSL2: \texttt{apt install tcsh} \\[1.mm]
Python, \texttt{pip}
  & distribution packages, or Miniforge
  & \texttt{python.org} installer, or Miniforge
  & \texttt{python.org} installer, or Miniforge \\[1.mm]
OpenTURNS
  & \texttt{pip install openturns}; \texttt{conda install openturns};
    Debian/Ubuntu and Fedora packages
  & \texttt{pip} or \texttt{conda}
  & \texttt{pip} or \texttt{conda} \\[1.mm]
JAX
  & \texttt{pip install -U jax} (CPU); \texttt{"jax[cuda13]"} for an NVIDIA GPU
  & \texttt{pip install -U jax} (Apple silicon, CPU)
  & \texttt{pip install -U jax} (CPU, experimental); GPU only through WSL2 \\
\end{tabular}
\end{gbox}

\noindent Three practical remarks.

\begin{itemize}
\item \textbf{Use one Python environment per project}, so that a library
  upgraded for another project does not break this one:
\begin{codebox}
\begin{verbatim}
$ python3 -m venv ~/venv/castem          # once
$ source ~/venv/castem/bin/activate      # bash/zsh; csh: activate.csh
$ pip install numpy scipy pandas matplotlib openturns jax
\end{verbatim}
\end{codebox}
  With Miniforge, \texttt{conda create -n castem -c conda-forge python
  openturns jax} does the same; OpenTURNS recommends Miniforge (conda-forge)
  rather than the Anaconda channels.
\item \textbf{On Windows, script from WSL2.} The shell drivers need a Unix
  shell; WSL2 (\texttt{wsl --install} in an administrator terminal) gives a
  Ubuntu system inside Windows, in which the Linux version of \texttt{Cast3M},
  the shells and all the Python libraries install as on Linux. The Python
  drivers also run in a native Windows Python, provided the \texttt{Cast3M}
  command they call is on the \texttt{PATH}: check its name in the shortcut
  the installer created.
\item \textbf{Test the drivers without \texttt{Cast3M}.} A small script that
  reads \texttt{calcul.dgibi} and writes the CSV file \texttt{Cast3M} would
  write, from a formula, is enough to develop a driver on a laptop or in a
  continuous-integration service where \texttt{Cast3M} is not installed.
  Name it \texttt{castem26}, put it first on the \texttt{PATH}, and replace
  it by the real code when the driver works.
\end{itemize}

\noindent The scripts of sections~\ref{sec:py:csh} to~\ref{sec:py:jax} were run
with \texttt{tcsh}~6.24, Python~3.13, OpenTURNS~1.27 and JAX~0.11, against
such a stand-in, not against a \texttt{Cast3M} installation.

\framepart{Summary}

\begin{gbox*}
\begin{itemize}
  \item \textbf{Export once, analyse outside.} Curves leave
  \texttt{Cast3M} as a CSV file with a header -- the \op{EXPORTCSV}
  procedure of section~\ref{sec:py:export} is the modern replacement for
  \op{"@excel1} -- and fields leave as VTK.
  \item \textbf{pandas and matplotlib for the curves}, and they give you
  vector figures that drop straight into \LaTeX{}.
  \item \textbf{pyvista for the fields}, which does things the built-in
  viewer cannot: sampling along an arbitrary path, integrating over the
  domain, comparing two different meshes on the same abscissa.
  \item \textbf{A parameter study is a template plus a loop.}
  \texttt{sed} substitutes the parameters into a \texttt{.dgibi} file, the
  shell runs each case \emph{in its own directory} -- \texttt{fort.3} and
  \texttt{fort.98} make that compulsory -- and \texttt{awk} reduces the
  results to one table. Python takes over wherever the analysis stops being
  column arithmetic.
  \item \textbf{\texttt{csh} for what you inherit, \texttt{bash} for what
  you write.} The loop is the same; only the syntax of loops, tests and
  redirections changes.
  \item \textbf{\texttt{Cast3M} as a function}: a template, a directory per
  run, a check for \texttt{ERREUR} and a cache make it usable by any driver --
  OpenTURNS for uncertainties and calibration, JAX when the cost around it is
  a program to differentiate, with the derivative of the black box declared by
  \texttt{custom\_vjp}.
  \item \textbf{The real gain is not the prettier picture}: it is that the
  whole study becomes a script. Rerun the computation, rerun the script, and
  every figure in your report updates.
\end{itemize}
\end{gbox*}

\framepart{Exercises}

\exercise{Replace \texttt{@excel1}} Take any computation from the
  previous chapters, export its results with \op{EXPORTCSV}, and reproduce in
  matplotlib a figure you had previously made with \op{dess}. Which of the
  two is faster to redo after the computation changes?

\exercise{Convergence study} Use the parametric pattern of
  section~\ref{sec:py:param} to run the plate with \op{nel} $=5,10,20,40,80$
  and plot the maximum von Mises stress against the number of degrees of
  freedom, on logarithmic axes. Estimate the convergence rate from the slope.

\exercise{The same sweep, twice} Do the convergence study of the
  previous exercise a second time, with a shell script instead of Python:
  a template plus \texttt{sed}, a loop over \op{nel}, and an \texttt{awk}
  one-liner to pull the maximum stress out of each \texttt{history.csv} into
  a single table. Which version was quicker to write, and which would you
  rather hand to someone else?

\exercise{Compare two meshers} Compute the same problem on the
  \texttt{Cast3M} mesh of chapter~\ref{chap:meshing} and on the
  \textsc{gmsh} mesh of section~\ref{sec:io:gmsh}, export both to VTK, and use
  \texttt{sample\_over\_line} to plot $\sigma_{yy}$ along the same physical
  path for both. This is a comparison the built-in \op{evol} cannot make,
  since the two meshes have different nodes.

\exercise{The sweep in C shell} Rewrite the \texttt{bash} script of your
  convergence study in \texttt{csh}, run both, and compare the two tables
  line by line with \texttt{diff}.

\exercise{Scatter of the stress concentration} With the OpenTURNS example of
  section~\ref{sec:py:openturns}, compute the mean and the 99\% quantile of the
  maximum stress for a radius of 5\% scatter. Then fit a polynomial chaos on
  20 runs (\texttt{FunctionalChaosAlgorithm}) and compare the same two
  numbers, computed on $10^5$ points of the metamodel.

\exercise{Checking a declared derivative} In the JAX example of
  section~\ref{sec:py:jax}, compute $e_1(h)=\abs{J(r+h)-J(r)-h\,J'(r)}$ for
  $h=10^{-1},\dots,10^{-4}$ and plot it on logarithmic axes. Where does the
  slope 2 stop, and why does the step of the difference in \texttt{\_df}
  decide it?
